\documentclass[11pt]{article}
\usepackage{fontspec,amsmath,amssymb,amsthm,hyperref}
\usepackage[margin=1in]{geometry}
\setlength{\emergencystretch}{3em}
\title{Marginal partitions, existence and the exact parameter target}
\author{LangevinLean mathematical exposition}
\date{}
\begin{document}\maketitle
\noindent Source module: \texttt{IPLA/MarginalConcentration.lean}. Compile this proof document with XeLaTeX or LuaLaTeX. The accompanying HTML pairs the mathematical argument with the exact Lean statements and records the scope of the formalization.


For differentiable jointly \(m\)-strongly convex \(U\), \(m>0\), this module proves finite positive latent partitions, continuity and vanishing parameter tails, existence of a maximizer, and the exact parameter projection of \(\Pi_N\). It establishes these facts before any logarithm, real conversion or concentration estimate is used.


\section{Pair coordinates split the squared norm}


The pairing map has parameter coordinates \(\theta_j\) and latent coordinates \(x_j\). Expanding the Euclidean squared norm over the disjoint index sets gives

\[\|\operatorname{pair}(\theta,x)\|^2=\|\theta\|^2+\|x\|^2.\]

Every output coordinate of pairing is a continuous product projection; the finite coordinate-to-Euclidean identification is continuous as well. Hence \((\theta,x)\mapsto\operatorname{pair}(\theta,x)\) is jointly continuous. These elementary identities justify the orthogonal Gaussian domination and later dominated convergence.




\textbf{Correspondence with the displayed declarations.}


\begin{itemize}

\item \texttt{norm\_pair\_sq}: The pair coordinates preserve the orthogonal decomposition into parameter and latent variables.
\item \texttt{continuous\_pair}: Joint continuity of the parameter/latent pairing.
\end{itemize}


\noindent Lean declarations: \texttt{pair\_parameter}, \texttt{pair\_particle}, \texttt{norm\_pair\_sq}, \texttt{continuous\_pair}.

\section{Strong convexity supplies uniform Gaussian domination}


The shared strong-convexity lower bound, obtained by bounding the linear term at zero, is

\[U(v)\ge U(0)-\frac{\|\nabla U(0)\|^2}{m}+\frac m4\|v\|^2.\]

Set \(C=\exp(\|\nabla U(0)\|^2/m-U(0))\) and \(a=m/4>0\). Exponentiation and the norm splitting yield

\[e^{-U(\theta,x)}\le Ce^{-a\|\theta\|^2}e^{-a\|x\|^2}\le Ce^{-a\|x\|^2}.\]

The final majorant is independent of \(\theta\) and integrable over the latent Euclidean space. Continuity supplies measurability, so domination proves integrability of every latent slice. No bounded parameter neighborhood is needed.




Strong convexity and the affine score term at zero give a Gaussian majorant after completing a square. Split the joint squared norm as \(\|\operatorname{pair}(\theta,x)\|^2=\|\theta\|^2+\|x\|^2\). Dropping the nonnegative parameter contribution yields an integrable Gaussian bound in \(x\), uniform in \(\theta\) when needed for local parameter limits. This proves latent integrability before converting the extended marginal density to its real integral. Strict positivity follows because the exponential weight is positive everywhere.




\textbf{Correspondence with the displayed declarations.}


\begin{itemize}

\item \texttt{marginal\_weight\_gaussian\_bound}: A Gaussian domination uniform in both the parameter and latent variable.
\item \texttt{marginal\_weight\_uniform\_bound}: An integrable bound for every parameter slice, with no local parameter restriction.
\item \texttt{integrable\_marginal\_weight}: Every latent partition is integrable from joint assumptions alone.
\end{itemize}


\noindent Lean declarations: \texttt{marginal\_weight\_gaussian\_bound}, \texttt{marginal\_weight\_uniform\_bound}, \texttt{integrable\_marginal\_weight}.

\section{The marginal is positive, finite and continuous}


Let \(M(\theta)=\int e^{-U(\theta,x)}dx\). Slice integrability and nonnegativity identify the extended marginal with \(\operatorname{ofReal}(M(\theta))\). This excludes infinity. The integrand is strictly positive everywhere, and Lebesgue measure is nonzero, so \(M(\theta)>0\) and the extended marginal is strictly positive as well. If \(\theta_n\to\theta\), the integrands converge pointwise by continuity of \(U\), and the preceding common Gaussian majorant permits dominated convergence. Thus \(M\) is continuous; composing with the continuous \(\operatorname{ofReal}\) map proves continuity of the extended marginal.




\textbf{Correspondence with the displayed declarations.}


\begin{itemize}

\item \texttt{marginalDensity\_eq\_ofReal\_integral}: The finite extended integral agrees with the ordinary real marginal.
\item \texttt{marginalDensity\_ne\_top}: The parameter marginal is finite, excluding infinite-likelihood maximizers.
\item \texttt{marginalDensity\_pos}: Every parameter has strictly positive marginal density.
\item \texttt{continuous\_marginal\_integral}: Continuity of the real-valued marginal follows by dominated convergence.
\item \texttt{continuous\_marginalDensity}: The actual extended-valued marginal is continuous as well.
\end{itemize}


\noindent Lean declarations: \texttt{marginalDensity\_eq\_ofReal\_integral}, \texttt{marginalDensity\_ne\_top}, \texttt{marginalDensity\_pos}, \texttt{continuous\_marginal\_integral}, \texttt{continuous\_marginalDensity}.

\section{Gaussian parameter tails force existence of a maximizer}


Integrating the stronger Gaussian domination gives

\[0<M(\theta)\le C\left(\int e^{-a\|x\|^2}dx\right)e^{-a\|\theta\|^2}.\]

The finite constant is independent of \(\theta\), and the right-hand side tends to zero as \(\theta\) leaves compact sets. The squeeze theorem proves \(M(\theta)\to0\) at infinity. Since \(M(0)>0\), outside a sufficiently large compact region the values are below \(M(0)\). Continuity then attains a global maximum within a compact region. Monotonicity of \(\operatorname{ofReal}\) transfers that maximum to the extended marginal predicate. Uniqueness is proved later from Prékopa strong convexity.




Integrating the stronger Gaussian bound while retaining the parameter term gives \(M(\theta)\le C e^{-c\|\theta\|^2}\), for finite \(C\) and \(c>0\). Dominated convergence proves continuity of \(M\), and this bound proves \(M(\theta)\to0\) as \(\|\theta\|\to\infty\). Since \(M\) is positive at any chosen point, outside a sufficiently large ball it is smaller than that value. Compact attainment on the ball proves existence of the actual global marginal maximizer.




\textbf{Correspondence with the displayed declarations.}


\begin{itemize}

\item \texttt{marginal\_integral\_gaussian\_bound}: Joint strong convexity gives a Gaussian parameter tail after integrating out the latent variable. The constant is finite and independent of \texttt{theta}.
\item \texttt{tendsto\_marginal\_integral\_cocompact}: The parameter marginal vanishes at infinity.
\item \texttt{exists\_marginalMaximizer}: Existence of the actual marginal maximizer is derived from the original joint assumptions. No maximizer or compactness premise is introduced.
\end{itemize}


\noindent Lean declarations: \texttt{marginal\_integral\_gaussian\_bound}, \texttt{tendsto\_marginal\_integral\_cocompact}, \texttt{exists\_marginalMaximizer}.

\section{Product integration and the particle-scaling Jacobian}


Tonelli induction over \(n\) proves the extended-integral identity

\[\int\prod_{i=1}^nf(x_i)\,d\mu^{\otimes n}=\left(\int f\,d\mu\right)^n.\]

The base case is the mass-one empty product; the successor case separates one coordinate and multiplies the remaining integral. This statement records divergence correctly because no real conversion is used. For a nonzero real dilation \(r\) in dimension \(d\), Haar measure scaling gives \(\int f(rx)dx=|r^d|^{-1}\int f(x)dx\). Applying these two identities at \(r=\sqrt N\), with \(N>0\), yields

\[\int\prod_i e^{-U(\theta,\sqrt N z_i)}\,dz_1\cdots dz_N=\left(N^{-d_x/2}M(\theta)\right)^N.\]

The formal statement retains the absolute-value inverse-power Jacobian, which is equal to the displayed factor.




For a fixed \(\theta\), the lifted weight factors as \(\prod_i e^{-U(\theta,\sqrt N z_i)}\). Tonelli factors the nonnegative product integral into a product of identical latent integrals. The substitution \(x_i=\sqrt N z_i\) contributes \(N^{-d_x/2}\) per particle, so the fiber integral is \(N^{-Nd_x/2}M(\theta)^N\). This Jacobian is finite and strictly positive when \(N>0\), and does not depend on \(\theta\).




\textbf{Correspondence with the displayed declarations.}


\begin{itemize}

\item \texttt{lintegral\_fin\_prod\_eq\_pow}: Tonelli factorization of identical nonnegative weights over finitely many independent latent coordinates. This extended-integral statement also records divergence correctly; no integrability premise is hidden.
\item \texttt{lintegral\_comp\_smul\_eq}: Change of variables for a nonnegative weight under scalar dilation.
\item \texttt{scaled\_particle\_partition}: Integrating the actual scaled particle potential at a fixed parameter gives the \texttt{N}th power of the original marginal, times the explicit scalar Jacobian. This is the factor that cancels in the normalized parameter law.
\end{itemize}


\noindent Lean declarations: \texttt{lintegral\_fin\_prod\_eq\_pow}, \texttt{lintegral\_comp\_smul\_eq}, \texttt{scaled\_particle\_partition}.

\section{A coordinate equivalence groups the state into \ensuremath{\theta} and particle blocks}


The configuration map \(C(\theta,(z_i)_i)\) inserts the parameter coordinates first and then all particle coordinates. Its coordinate evaluation theorems follow directly from insertion. The inverse is \(z\mapsto(p_\theta z,(\operatorname{block}_iz)_i)\); checking each coordinate proves both inverse identities and linearity, hence a finite-dimensional continuous linear equivalence. The evaluation theorem identifies this equivalence with the configuration definition, while \(p_\theta C(\theta,z)=\theta\) and \(A_iC(\theta,z)=\operatorname{pair}(\theta,\sqrt N z_i)\). These identities give the exact arguments of the fiber integrals.




\textbf{Correspondence with the displayed declarations.}


\begin{itemize}

\item \texttt{configuration}: Assemble the scaled state from its parameter and particle blocks.
\item \texttt{configurationCLE}: Grouping coordinates is a constructed continuous linear equivalence.
\end{itemize}


\noindent Lean declarations: \texttt{configuration}, \texttt{configuration\_parameter}, \texttt{configuration\_particle}, \texttt{configurationCLE}, \texttt{configurationCLE\_apply}, \texttt{parameter\_configuration}, \texttt{joint\_configuration}.

\section{The weight factors and its parameter projection integrates each fiber}


In grouped coordinates, exponentiation converts the sum in \(V_N\) into a product:

\[e^{-V_N(C(\theta,z))}=\prod_i e^{-U(\theta,\sqrt N z_i)}.\]

Two measure identities justify using this pointwise formula. First, for a measurable equivalence \(e\), pushing \(\mu\) with weight \(w\circ e\) through \(e\) gives \((e_\#\mu)\) with weight \(w\); prove this by evaluating any measurable set and the pushforward integral formula. Second, projecting \(w(\theta,z)d\mu(\theta)d\nu(z)\) to its first coordinate gives density \(\int w(\theta,z)d\nu(z)\) relative to \(\mu\), by Tonelli on the indicator of an arbitrary measurable parameter set.




\textbf{Correspondence with the displayed declarations.}


\begin{itemize}

\item \texttt{weight\_configuration}: In grouped coordinates the actual Gibbs weight factors over particles.
\item \texttt{map\_withDensity\_comp\_equiv}: A density commutes with a measurable coordinate equivalence.
\item \texttt{map\_fst\_withDensity}: Tonelli's theorem identifies the first-coordinate density of a weighted product measure.
\end{itemize}


\noindent Lean declarations: \texttt{weight\_configuration}, \texttt{map\_withDensity\_comp\_equiv}, \texttt{map\_fst\_withDensity}.

\section{All finite positive volume constants cancel on normalization}


A continuous linear equivalence carries Euclidean volume to a Haar measure, hence to \(c\) times the target volume for a finite positive \(c\). This proves the required coordinate-volume relation without assuming its scalar is one. For any measure \(\mu\) and \(0<c<\infty\), extended-real inverse multiplication gives

\[\frac{c\mu}{(c\mu)(E)}=\frac{\mu}{\mu(E)}.\]

Thus both the coordinate-grouping scalar and the particle-dilation Jacobian cancel in a normalized law. The formal normalization lemma proves this algebra for arbitrary \(\mu\), keeping the nonzero and finite conditions on \(c\) explicit.




\textbf{Correspondence with the displayed declarations.}


\begin{itemize}

\item \texttt{normalize\_smul\_eq}: Normalization is unchanged by multiplication by any finite positive constant. This is why grouping coordinates and the particle Jacobian cancel.
\item \texttt{configuration\_map\_volume}: The finite-dimensional coordinate grouping changes volume by a constructed finite positive constant. Its value is irrelevant after normalization.
\end{itemize}


\noindent Lean declarations: \texttt{normalize\_smul\_eq}, \texttt{configuration\_map\_volume}.

\section{The actual parameter marginal is the powered likelihood target}


Define the normalized power target \(\pi_N^\theta(d\theta)=Z_\theta^{-1}M(\theta)^N d\theta\). To identify it, let \(\mu=e^{-V_N}dz\) and \(\rho=M(\theta)^Nd\theta\). Grouping coordinates changes volume by \(c\), and the fiber calculation contributes \(J^N\), \(J=N^{-d_x/2}\). The preceding pushforward identities prove \(c\,p_{\theta\#}\mu=J^N\rho\). Both scalars are finite and positive. Normalize this equality, cancel both scalars, and use preservation of total mass under measurable projection to conclude

\[p_{\theta\#}\Pi_N=\pi_N^\theta.\]

Every change of variables and normalization factor in this identity is proved; no formula for the marginal is assumed.




The configuration equivalence groups the orthogonal state coordinates into the parameter block and all latent blocks. The volume theorem tracks the finite-dimensional coordinate-change constant, while the density pushforward lemmas integrate each parameter fiber. The preceding factorization leaves a density proportional to \(M(\theta)^N\), with the scaling and coordinate-change factors independent of \(\theta\). Normalization cancels those positive finite factors exactly, yielding \[(p_\theta)_\#\Pi_N(d\theta)=\frac{M(\theta)^N\,d\theta}{\int M(u)^Ndu}.\] The final target probability and moment lemmas are consequences of this identity and the already proved lifted target facts.




\textbf{Correspondence with the displayed declarations.}


\begin{itemize}

\item \texttt{marginalTarget}: The parameter law with density proportional to the \texttt{N}th power of the original latent marginal. This names the statistical target without assuming anything about its concentration.
\end{itemize}


\noindent Lean declarations: \texttt{marginalTarget}, \texttt{gibbsTarget\_parameter\_marginal}.

\section{Normalize the power target and guard its moments}


The established projection identity and the probability property of \(\Pi_N\) prove that \(\pi_N^\theta\) is a probability measure. Its unnormalized density \(M^N\) is everywhere positive, so its extended partition is strictly positive. If that partition were infinite, the definition of the normalized target would be the zero measure, contradicting mass one; thus it is also finite. Finally the lifted Gibbs law has a finite second moment and \(p_\theta\) is a bounded linear map. The pushforward integral identity therefore gives \(\theta\mapsto\theta\in L^2(\pi_N^\theta)\). This supplies finite marginal moments before the later sharp bound is applied.




\textbf{Correspondence with the displayed declarations.}


\begin{itemize}

\item \texttt{marginalTarget\_isProbability}: The marginal power density normalizes to a probability measure as a consequence of its proved equality with the actual Gibbs projection.
\item \texttt{marginalPartition\_pos}: The normalizer of the powered marginal density is strictly positive.
\item \texttt{marginalPartition\_ne\_top}: The powered marginal partition is finite. Its proof uses the previously established probability identity, so divergence cannot disappear in a real conversion or an infinite normalizing constant.
\item \texttt{marginalTarget\_memLp}: The parameter marginal has finite second moment, without a separate marginal moment or concentration assumption.
\end{itemize}


\noindent Lean declarations: \texttt{marginalTarget\_isProbability}, \texttt{marginalPartition\_pos}, \texttt{marginalPartition\_ne\_top}, \texttt{marginalTarget\_memLp}.

\end{document}
